\section{Mechanism of Forking}
\label{sec:mechanism}

The approximate power law in Figure~\ref{fig:frequency} is our first clue: why should context frequency organize the gap so simply?

\subsection{From the frequency law to novel continuations}
\label{sec:coverage}
A rare context has had few chances to reveal its possible continuations. Call a validation target \emph{novel} when its context occurs in training but was never followed there by that target, and let $f(F)$ be the fraction of validation positions at context count $F\ge1$ whose continuation is novel. This fraction closely tracks the frequency dependence of the gap: the dashed curves in Figure~\ref{fig:frequency} are the directly measured $f$, rescaled and shifted,
\begin{equation}
\gap_t(F)\approx a_t f(F)+b_t,
\label{eq:coverage-main}
\end{equation}
with one scale and one offset per branch at training step $t$. Good--Turing estimation connects this novel fraction to singleton counts in the corpus; heavy-tailed conditional distributions provide a reason for its slow, approximately power-law decay. The detailed derivation and the terms beyond simple proportionality are in Appendix~\ref{app:math}. Thus the frequency trend points to a property of the text: rare contexts leave more plausible continuations unseen. It remains to explain why those continuations are harmed by replay.

\subsection{Crowding-out under replay}
\label{sec:sharpening}
A finite corpus does not fully represent the continuation distribution of each context: a context whose population distribution assigns equal probability to $C$ and $D$ may appear three times with $C$ and once with $D$ in training, so fitting the sample underpredicts $D$ on fresh data. The replay trajectory adds a further clue. After one complete epoch, the corpus and its novel fraction are fixed, yet the gap continues to grow, by comparable amounts over the first few epochs (Figure~\ref{fig:staircase}; Appendix~\ref{app:replay-clue}). The missing coverage alone does not specify this continuing damage. We hypothesize that each repetition reinforces the seen continuations and progressively suppresses the alternatives.

This would make forking ordinary overfitting confined to very small problems. For a context $c_i$, call the next-token logits computed from it its \emph{context subspace}, and split the vocabulary into the continuations of $c_i$ seen in training, those that appear only in validation, and the rest (Appendix~\ref{app:crowding-derivation}). If training shows a single continuation $C$, the loss in that subspace is $-\ln p_C$; an idealized step directly in logit coordinates raises the logit of $C$ by $\eta(1-p_C)$ and, through the softmax, lowers every other logit by $\eta p_v$. Mass flows from the unseen continuations onto the seen one, so the seen-to-unseen odds grow by a factor above one at every epoch with nothing to stop them while $p_C<1$ (Appendix Equation~\eqref{eq:zt}). We call this \emph{crowding-out}; it provides an idealized account of the growth in the novel-continuation penalty with replay.


\begin{figure}[!ht]
\centering
\includegraphics[width=\linewidth]{figures/fig_replay_transfer_schematic.pdf}
\caption{Local sharpening and cross-context recovery, illustrated schematically. Dashed marks show prior heights; darker extensions show gains. \emph{Others} denotes the mean probability of the top 20 remaining tokens. Measured distributions for \emph{there is a} appear in Figure~\ref{fig:actual-logits-examples}(e,f).}
\label{fig:replay-transfer}
\end{figure}

Figure~\ref{fig:replay-transfer} illustrates the proposed mechanism: replay sharpens a local distribution, while updates from related contexts can support alternative continuations. Its heights are illustrative. Appendix~\ref{app:sharpening} presents the measured context distributions, including \emph{petting your}, together with aggregate seen/novel losses and predictive entropy (Figure~\ref{fig:subspace-ce}). The aggregate measurements show worsening novel-continuation accuracy in the both-table model while the no-gram control improves; individual contexts need not sharpen monotonically. Crowding-out also offers an explanation for continued forking after the tables are frozen (Section~\ref{sec:interventions}): the context-specific representation remains available, and backbone updates can continue to favor the observed continuations.

This leaves one question: why does the same replay improve the no-gram model on novel continuations, while the $n$-gram model grows less accurate? The next experiment shows that the table lets sharpening proceed independently within each context.

\FloatBarrier
\subsection{Context subspaces are independent}
\label{sec:context-locality}
Section~\ref{sec:sharpening} treats each context subspace as evolving on its own. We test this by ablation: from the start of epoch 2 or 3, we exclude the next-token loss at a random 2\% of training trigram contexts, keeping their inputs and all parameters trainable, while a disjoint 2\% serve as unmasked peers (Appendix~\ref{app:context-locality}). Once the mask starts, the masked cohort's mean gap stops growing (Figure~\ref{fig:subspace-gap}), whereas the peers' and normal training's jump again at every replay boundary. In paired one-epoch runs, masking reduces the epoch-2 gap increase of the selected contexts from $1.018$ to $0.007$ nats, while the peers' increase stays at $0.997$--$0.998$. Gap growth is generated inside each context's own subspace, not imported from the other 98\% of contexts that keep training.

\begin{figure}[!t]
\centering
\begin{subfigure}[b]{0.36\linewidth}
\centering
\includegraphics[width=\linewidth]{figures/fig_mechanism_a.pdf}
\caption{}
\label{fig:subspace-mass}
\end{subfigure}\hfill
\begin{subfigure}[b]{0.34\linewidth}
\centering
\includegraphics[width=\linewidth]{figures/fig_mechanism_b.pdf}
\caption{}
\label{fig:subspace-gap}
\end{subfigure}\hfill
\begin{subfigure}[b]{0.28\linewidth}
\centering
\includegraphics[width=\linewidth]{figures/fig_mechanism_c.pdf}
\caption{}
\label{fig:subspace-ce}
\end{subfigure}
\caption{Context-subspace evidence for replay-driven sharpening (seed 42). (a) Mean probability mass on train-seen continuations (top) and on validation-only continuations (bottom) over four epochs, for masked and unmasked context cohorts in the $n$-gram and no-gram models. (b) Mean fixed-probe gap of the masked cohort, of its unmasked peers, and of normal training; dotted lines mark epoch boundaries. (c) Validation cross-entropy at the five epoch-end checkpoints for train-seen bigram contexts, split by whether the true continuation appeared in training.}
\label{fig:subspace-evidence}
\end{figure}

The mass on train-seen continuations shows the same partition (Figure~\ref{fig:subspace-mass}). In the $n$-gram model, unmasked contexts raise it from $0.43$ to $0.69$ by epoch 4, while masked contexts end at $0.40$. Without the $n$-gram module, masked contexts still gain mass ($0.22\to0.39$, versus $0.41$ unmasked), and replay also raises the validation-only mass ($0.013\to0.06$). The table therefore partitions optimization into weakly interacting context subspaces, inside which replay-driven sharpening proceeds unchecked by other contexts. The shared backbone still moves individual masked contexts, but it shifts their training and validation losses together and leaves their average gap flat.

\FloatBarrier
